\documentclass[11pt]{article}
\usepackage[margin=1in]{geometry}
\usepackage{amsmath,amssymb}
\newcommand{\finalanswer}[1]{\par\medskip\noindent\fbox{\parbox{0.94\linewidth}{\textbf{Answer:} #1}}}
\title{STAT 412 Homework 6 --- Question 3}
\author{}\date{}
\begin{document}\maketitle
\section*{Problem}
For $n=59$, $\bar x=82.38^\circ$F, and known $\sigma=13.78^\circ$F, construct 90\% and 95\% confidence intervals for $\mu$.
\section*{Work}
Since $\sigma$ is known, use
\[
\bar x\pm z_{\alpha/2}\frac{\sigma}{\sqrt n}.
\]
For 90\%, $z_{\alpha/2}=1.645$:
\[
E=1.645\frac{13.78}{\sqrt{59}}\approx2.95,
\]
so
\[
82.38\pm2.95=(79.43,85.33).
\]
For 95\%, $z_{\alpha/2}=1.960$:
\[
E=1.960\frac{13.78}{\sqrt{59}}\approx3.52,
\]
so
\[
82.38\pm3.52=(78.86,85.90).
\]
The 95\% interval is wider because the higher confidence level uses the larger critical value. The correct interpretation is choice B: you can be 90\% confident that the population mean record high temperature is between the bounds of the 90\% confidence interval, and 95\% confident for the 95\% interval.
\finalanswer{90\% CI $(79.43,85.33)$; 95\% CI $(78.86,85.90)$. The 95\% interval is wider. Interpretation choice B.}
\end{document}
